\subsection{The Atomic Wiring Algorithm}

Let a Blueprint be a finite partial mapping
\[
BP : K \rightharpoonup V
\]
where $V$ may denote:
\begin{itemize}
    \item Static values,
    \item Reactive signals,
    \item Nested blueprint objects,
    \item Lifecycle hooks.
\end{itemize}

The reification process constructs a \textbf{Persistent Shell}
\[
E = (N, \mathcal{F})
\]
where:
\begin{itemize}
    \item $N$ is the underlying DOM node,
    \item $\mathcal{F}$ is the set of atomic reactive bindings.
\end{itemize}

\subsubsection*{Deterministic Atomic Wiring}

The construction phase performs a single linear scan over $BP$.
No recursive structural reconciliation is required after initialization.

\begin{lstlisting}[caption={Deterministic Atomic Wiring Algorithm}]
Algorithm: Reify(BP)

Input: Blueprint BP
Output: Persistent Shell E = (N, F)

1. N = CreateElement(BP.tag)
2. F = {}

3. for each (k, v) in BP do
4.     if IsReactive(v) then
5.         define f_k(s) = Update(N, k, s)
6.         RegisterBinding(v, f_k)
7.         F = F union {f_k}

8.     else if IsNestedBlueprint(v) then
9.         child = Reify(v)
10.        AttachChild(N, child)

11.    else if IsLifecycleHook(k) then
12.        RegisterHook(N, k, v)

13.    else
14.        ApplyStatic(N, k, v)

15. return (N, F)
\end{lstlisting}

\paragraph{Invariant.}
The dependency graph between reactive signals and DOM attributes
is constructed exactly once during the reification phase.
After this stage, the structural topology of the UI remains fixed.
State updates do not trigger structural traversal.

\paragraph{Time Complexity.}

Let $|BP| = n$.  
The reification procedure performs a single linear pass:

\[
T_{\text{reify}}(n) = O(n)
\]

Let $d$ denote the maximum number of attributes dependent on a given signal.

For a signal update:

\[
T_{\text{update}} = O(d)
\]

Under bounded dependency ($d \leq c$ for constant $c$):

\[
T_{\text{update}} = O(1)
\]

\paragraph{Deterministic Update Path.}
Unlike tree-diffing reconciliation strategies with complexity $O(N)$,
the CARM model guarantees that only directly affected attributes are updated.
This yields a constant-time atomic mutation path independent of global tree size.
